% W4i47_Observables.tex
% DFD 17-Agent Research Programme --- Iteration 47 (i47)
% Agents 8 (Laboratory), 9 (Particle Physics), 10 (Astrophysics), 11 (CMB), 12 (LSS), 13 (Dark Sector), 14 (SymBoltz/Numerics)
% Theme: Phase 0A EXECUTION, birefringence npi escalation, neutrino mass tension resolution paths, Euclid 85 days
% Date: 2026-03-23

\documentclass[11pt,a4paper]{article}
\usepackage[utf8]{inputenc}
\usepackage{amsmath,amssymb,amsfonts,amsthm}
\usepackage{hyperref}
\usepackage{booktabs}
\usepackage{longtable}
\usepackage{geometry}
\usepackage{xcolor}
\usepackage{tcolorbox}
\geometry{margin=2.5cm}

\newtheorem{theorem}{Theorem}
\newtheorem{proposition}{Proposition}
\newtheorem{lemma}{Lemma}
\newtheorem{corollary}{Corollary}
\newtheorem{definition}{Definition}
\newtheorem{remark}{Remark}

\definecolor{dfdgreen}{RGB}{0,128,0}
\definecolor{dfdyellow}{RGB}{200,180,0}
\definecolor{dfdred}{RGB}{200,0,0}
\definecolor{dfdblue}{RGB}{0,50,180}

\newcommand{\GREEN}{\textcolor{dfdgreen}{\textbf{GREEN}}}
\newcommand{\YELLOW}{\textcolor{dfdyellow}{\textbf{YELLOW}}}
\newcommand{\RED}{\textcolor{dfdred}{\textbf{RED}}}
\newcommand{\NEW}{\textcolor{dfdblue}{\textbf{NEW}}}
\newcommand{\RESOLVED}{\textcolor{dfdgreen}{\textbf{RESOLVED}}}
\newcommand{\Tr}{\operatorname{Tr}}
\newcommand{\CP}{\mathbb{C}P}

\title{DFD i47: Observables --- Laboratory, Particle, Astro, CMB, LSS, Dark Sector, Numerics\\[6pt]
\large Agents 8, 9, 10, 11, 12, 13, 14\\[6pt]
\normalsize Iteration 16/20 of Quantum Gravity Cycle (i32--i51)}
\author{DFD 17-Agent Research Programme}
\date{2026-03-23}

\begin{document}
\maketitle
\tableofcontents
\newpage

%=============================================================================
\part{Agent 8: Laboratory Experiments}
\label{part:agent8}
%=============================================================================

\section{Mandate}

Update Th-229 nuclear clock progress. SRF cavities. SQMS. Track $\dot{\alpha}/\alpha$ timeline.

\section{Th-229 Nuclear Clock: Programme Status Consolidated}

\begin{tcolorbox}[colback=green!5!white, colframe=green!60!black,
  title={Th-229 Programme: Three Pillars Established}]

The Th-229 nuclear clock programme now rests on three published pillars:
\begin{enumerate}
\item \textbf{Frequency reproducibility}: 220~Hz ($1.1 \times 10^{-13}$) at 195~K, demonstrated for 7 months (Nature, January 2026).
\item \textbf{Sensitivity coefficient}: $K_\alpha = 5900(2300)$, confirming three-orders-of-magnitude enhancement over atomic clocks (Nature Communications, 2025, now peer-reviewed).
\item \textbf{Systematic floor}: $10^{-18}$ fractional uncertainty achievable at optimal temperature ($196 \pm 5$~K).
\end{enumerate}

\textbf{Path to DFD test}: At $K_\alpha \approx 5900$ and $10^{-18}$ systematic floor, the nuclear clock can probe $\dot{\alpha}/\alpha \sim 10^{-19}$/yr with $\sim 10^4$ seconds of integration per measurement. Multi-crystal arrays (solid-state advantage) accelerate this. Timeline: $\sim$2030 for first physics-level constraint.

\textbf{Status}: \GREEN. On track. No new data since i46 but programme trajectory confirmed.
\end{tcolorbox}

\section{SRF and SQMS: Unchanged}

\begin{itemize}
\item Dark SRF refined bounds (PRL, March 2026) stand.
\item SQMS 2.0 (\$125M) continues.
\item Nb$_3$Sn cavity quality at elevated temperature: ongoing data collection.
\end{itemize}

\section{Agent 8 Assessment: Grade A-}

Unchanged from i46. No new laboratory results but programme on track.

\section{New Literature (Agent 8)}

\begin{enumerate}
\item \textbf{Nature (January 2026)} --- Th-229 frequency reproducibility. (continued; status: published)
\item \textbf{Nature Communications (2025)} --- $K_\alpha = 5900(2300)$. (continued; now peer-reviewed)
\item \textbf{Dark SRF} (PRL, March 2026) --- Refined bounds. (continued)
\item \textbf{SQMS 2.0} --- \$125M renewal. (continued)
\item \textbf{Dark matter via Th-229} (July 2025) --- Detection method. (continued)
\end{enumerate}


%=============================================================================
\part{Agent 9: Particle Physics}
\label{part:agent9}
%=============================================================================

\section{Mandate}

LHC Run 3 final months. Moriond 2026 follow-up. HL-LHC preparations.

\section{Moriond 2026: Full Digest --- No DFD Surprises}

\begin{tcolorbox}[colback=green!5!white, colframe=green!60!black,
  title={Moriond 2026 Complete: 45 Analyses, All SM-Consistent}]

The full Moriond 2026 results (15--22 March) are now digested:
\begin{itemize}
\item $H \to \mu\mu$: $3.4\sigma$ evidence (ATLAS). SM-consistent.
\item $H \to Z\gamma$: $2.5\sigma$ combined Run 2 + Run 3. SM one-loop.
\item $HH \to bb\gamma\gamma$: Record limits on Higgs self-coupling. Enhanced by combining Run 2 + early Run 3.
\item SUSY: Record-setting exclusion limits across all channels. No signals.
\item $\tau \to 3\mu$: Lepton-flavour violation search. No signal.
\item $t\bar{t}$ resonances: No excess.
\end{itemize}

\textbf{DFD assessment}: Complete consistency with SM. The continued absence of BSM particles at the LHC is exactly what DFD predicts. The Higgs self-coupling limits will tighten dramatically at HL-LHC (Run 4, June 2030). \GREEN.
\end{tcolorbox}

\section{LHC Run 3: Final 3 Months}

\begin{itemize}
\item Run 3 ends late June 2026 (confirmed).
\item LS3 begins July 2026.
\item HL-LHC full-scale MQXFB quadrupole magnet tests: $\sim$2/3 complete.
\item HL-LHC Run 4 start: June 2030.
\end{itemize}

\section{Agent 9 Assessment: Grade B+}

Unchanged from i46.

\section{New Literature (Agent 9)}

\begin{enumerate}
\item \textbf{ATLAS Moriond 2026 summary} --- 45 analyses, no BSM. (continued; full digest)
\item \textbf{ATLAS $HH \to bb\gamma\gamma$} (March 2026) --- Record Higgs self-coupling limits. \NEW.
\item \textbf{ATLAS $\tau \to 3\mu$} (March 2026) --- LFV search. \NEW.
\item \textbf{HL-LHC MQXFB tests} (Fermilab, 2026) --- (continued)
\item \textbf{Higgs mass ATLAS} --- $125.11 \pm 0.11$ GeV. (continued)
\item \textbf{Higgs mass CMS} --- $125.35 \pm 0.15$ GeV. (continued)
\end{enumerate}


%=============================================================================
\part{Agent 10: Astrophysics}
\label{part:agent10}
%=============================================================================

\section{Mandate}

Wide binary update. MOND tests. Gaia DR4 timeline.

\section{\NEW: Wide Binary Tests --- Quality Framework Debate Continues}

\begin{tcolorbox}[colback=blue!5!white, colframe=blue!60!black,
  title={Wide Binary Gravity Tests: Methodological Convergence Toward GR}]

The wide binary gravity test literature continues to develop:
\begin{itemize}
\item The quality framework paper (MNRAS 547, 2026) emphasizes that undetected close binaries, chance alignments, and projection effects can \textbf{mimic} MOND-like signals.
\item Application of the quality framework (arXiv:2602.24035) shows no evidence for MOND.
\item A new study (arXiv:2504.07569, published April 2025) tests GR vs MOND with realistic triple modelling, finding results consistent with GR.
\item However, some authors maintain that five studies spanning 2023--2024 show MOND-consistent anomalies at low accelerations. The debate hinges on contamination modelling and sample selection.
\end{itemize}

\textbf{DFD assessment}: DFD predicts no MOND-like departure at any separation. The methodological convergence toward GR (with improving quality frameworks) supports DFD. Gaia DR4 (December 2026) with orbital solutions will be the definitive test. \GREEN{} for DFD.
\end{tcolorbox}

\section{Gaia DR4: December 2026 --- Unchanged}

ESA confirms Gaia DR4 no earlier than December 2026. DR4 will include epoch-level astrometry enabling self-consistent Keplerian orbital solutions. This is the definitive wide binary dataset.

\section{Agent 10 Assessment: Grade B+}

Unchanged from i46.

\section{New Literature (Agent 10)}

\begin{enumerate}
\item \textbf{arXiv:2504.07569} (April 2025) --- Wide binaries from Gaia DR3: GR vs MOND with triple modelling. \NEW.
\item \textbf{MNRAS 547 (2026)} --- Quality framework for wide binary tests. (continued)
\item \textbf{arXiv:2602.24035} (February 2026) --- No evidence for MOND. (continued)
\item \textbf{ESA Gaia DR4 scenario} --- December 2026. (continued)
\item \textbf{arXiv:2510.17746} --- ML astrometric exoplanet orbits. (continued)
\end{enumerate}


%=============================================================================
\part{Agent 11: CMB and Birefringence}
\label{part:agent11}
%=============================================================================

\section{Mandate}

Birefringence monitoring. $A_L$. CMB anomalies. SO calibration.

\section{\NEW: Birefringence --- $n\pi$ Phase Ambiguity Changes the Landscape}

\begin{tcolorbox}[colback=red!5!white, colframe=red!60!black,
  title={\NEW: PRL 2026 --- $n\pi$ Phase Ambiguity of Cosmic Birefringence}]

\textbf{The problem}: The birefringence angle $\beta \approx 0.3^\circ$ from Planck PR4 EB correlation has a degeneracy: $\beta + n\pi$ ($n = 0, 1, 2, \ldots$) all produce the same EB signal at leading order. A new PRL paper (Naokawa \& Namikawa, January 2026) shows that the EB spectrum \emph{shape} contains information to break this degeneracy.

\textbf{Key finding}: The true rotation angle may be \textbf{larger than $0.3^\circ$}. The shape analysis suggests $n > 0$ cannot be excluded.

\textbf{Impact on DFD accommodation options}:

\begin{center}
\begin{tabular}{lp{4cm}p{4cm}}
\toprule
\textbf{Option} & \textbf{If $\beta \approx 0.3^\circ$} & \textbf{If $\beta \approx 180.3^\circ$} \\
\midrule
1. Independent source & Easily accommodated & Requires very strong ALP coupling \\
2. $\psi F \tilde{F}$ & Small coupling needed & \textbf{Large coupling --- constrained} \\
3. CS evolution & Mild CS modification & \textbf{Major CS modification --- may conflict with $\alpha$} \\
\bottomrule
\end{tabular}
\end{center}

\textbf{Assessment}: Risk upgraded from MEDIUM-HIGH to \textbf{HIGH} for the scenario $n > 0$. However, the $n = 0$ (small angle) scenario remains the simplest interpretation. Simons Observatory ($\sim$2027) will be decisive.
\end{tcolorbox}

\section{\NEW: Galaxy Polarization --- Alternative Birefringence Probe}

\begin{tcolorbox}[colback=blue!5!white, colframe=blue!60!black,
  title={PRL 134, 161001 (2025): New Birefringence Probe Using Galaxy Polarization}]

A new method uses the correlation between integrated radio polarization direction of spiral galaxies and their apparent optical shapes to measure cosmic birefringence independently of CMB systematics. Per-galaxy uncertainty: $5$--$15^\circ$. SKA continuum surveys + optical shape catalogues promise comparable noise power spectrum to CMB-S4 with \textbf{different systematics}.

\textbf{DFD significance}: This provides an independent check on the CMB birefringence signal. If the galaxy method confirms birefringence, the signal is astrophysical (not a CMB calibration artifact). If it does not, the CMB signal may be systematic. Either way, DFD benefits from a cleaner measurement.
\end{tcolorbox}

\section{\NEW: Birefringence--Neutrino Mass Connection (Namikawa PRL)}

The PRL paper (arXiv:2506.22999, Namikawa et al.) proposing that birefringence resolves the negative $\Sigma m_\nu$ tension (discussed in detail in Agent 1) has direct implications for Agent 11's CMB analysis. If this mechanism is correct, the low-$\ell$ EE spectrum shape is modified by birefringence during reionization, affecting the extraction of $\tau$ and hence $\Sigma m_\nu$.

\textbf{Implication for $A_L$}: The $A_L$ extraction could also be affected if birefringence modifies the lensing-like smearing of peaks. However, $A_L$ is primarily a high-$\ell$ effect while birefringence acts at all $\ell$. The cross-talk is expected to be small but should be checked in Phase 0B.

\section{$A_L$: Still Waiting for Phase 0A}

$A_L = 1.025 \pm 0.017$ (Qu et al., PRL 136, 2026) remains the latest measurement. DFD prediction: $1.00$--$1.05$. \YELLOW{} pending Phase 0A execution.

\section{Agent 11 Assessment: Grade B+}

Unchanged from i46. The $n\pi$ phase ambiguity and galaxy polarization probe are significant new literature. The birefringence risk landscape is now more complex.

\section{New Literature (Agent 11)}

\begin{enumerate}
\item \textbf{PRL (2026)} --- $n\pi$ phase ambiguity of cosmic birefringence. \NEW.
\item \textbf{PRL 134, 161001 (2025)} --- Galaxy polarization birefringence probe. \NEW.
\item \textbf{PRL (2025)} --- Resolving negative $\Sigma m_\nu$ with birefringence (arXiv:2506.22999). \NEW.
\item \textbf{arXiv:2504.13154} (June 2025) --- Anisotropic birefringence null. (continued)
\item \textbf{arXiv:2602.12019} (February 2026) --- Simulation-based birefringence inference. (continued)
\item \textbf{arXiv:2512.19102} (December 2025) --- SO SAT polarization calibration. (continued)
\item \textbf{Qu et al.\ (PRL 136, 2026)} --- $A_L = 1.025 \pm 0.017$. (continued)
\end{enumerate}


%=============================================================================
\part{Agent 12: Large-Scale Structure and Euclid}
\label{part:agent12}
%=============================================================================

\section{Mandate}

Euclid pre-registration. DR1 preparations. DESI interplay.

\section{Euclid Countdown: 88 Days (20 June 2026)}

\begin{tcolorbox}[colback=yellow!5!white, colframe=yellow!60!black,
  title={Euclid Pre-Registration: 88 Days --- Phase 0A Still Required}]

The Euclid pre-registration deadline remains 20 June 2026 (88 days from i47).

\textbf{Phase 0A dependency}: The pre-registration requires at minimum a qualitative DFD prediction for Euclid observables ($D_L$ distance bias, lensing amplitude, BAO distance ratios). Phase 0A would provide quantitative numbers. Without Phase 0A, only a qualitative/parametric pre-registration is possible.

\textbf{Qualitative pre-registration draft}: Can proceed immediately with:
\begin{enumerate}
\item Statement of DFD's $\psi$-field distance-bias mechanism.
\item Parametric prediction: $D_L^{\text{DFD}} = D_L^{\Lambda\text{CDM}} \cdot e^{\Delta\psi(z)}$.
\item Range of allowed $\Delta\psi(z)$ from DFD field equations (without SymBoltz computation).
\item Fisher forecast for Euclid sensitivity to $\Delta\psi$ parameters.
\end{enumerate}

\textbf{Recommendation}: Begin qualitative pre-registration \textbf{immediately}, independent of Phase 0A. Phase 0A numbers can be added as a supplement before the deadline.
\end{tcolorbox}

\section{Euclid DR1: 21 October 2026}

\begin{itemize}
\item DR1 confirmed for 21 October 2026.
\item Higher-order WL statistics (Euclid Prep.\ LXXXV, A\&A, March 2026) outperform two-point by $2.5\times$ on $w_0$.
\item Quick Data Release 1 (QDR1) already published (March 2025): 30 TB data, deep fields.
\item ERO weak lensing pipeline validated (arXiv:2507.07629).
\end{itemize}

\section{\NEW: DESI DR2 Dynamical Dark Energy Robustness Debate}

\begin{tcolorbox}[colback=blue!5!white, colframe=blue!60!black,
  title={\NEW: DESI DR2 Dynamical DE --- Is the Evidence Robust?}]

The DESI DR2 $w_0 w_a$CDM preference ($2.8$--$4.2\sigma$) has generated a vigorous debate:

\textbf{For robustness}: Multiple analyses confirm the Quintom-B regime ($w_0 > -1$, $w_a < 0$) regardless of DE parameterization (arXiv:2511.22512, published March 2026; Nature Astronomy, 2025).

\textbf{Against robustness}: arXiv:2504.15222 (April 2025) argues that the combinations are problematic due to tensions among CMB, BAO, and SN datasets, concluding the claim is not robust. Dataset-independent analysis shows the real parameter space remains elusive.

\textbf{DFD assessment}: DFD is agnostic on $w_0 w_a$CDM vs $\Lambda$CDM at the background level. DFD's $\psi$-field produces a \emph{distance bias} that can mimic dynamical DE in BAO observations. If the DESI signal is real, it could be:
\begin{enumerate}
\item Genuine dynamical DE (not predicted by DFD at this level).
\item A $\psi$-field distance bias misinterpreted as dynamical DE.
\item A statistical fluctuation or systematic.
\end{enumerate}

Option 2 is testable: Phase 0A would produce the DFD prediction for BAO distance ratios. If DFD's $\psi$-field distance bias matches the DESI signal, this would be a \textbf{major success}. This adds urgency to Phase 0A execution.
\end{tcolorbox}

\section{Agent 12 Assessment: Grade B+}

Unchanged from i46. The DESI robustness debate adds context. Euclid pre-registration recommendation to proceed qualitatively is new.

\section{New Literature (Agent 12)}

\begin{enumerate}
\item \textbf{arXiv:2504.15222} (April 2025) --- Did DESI DR2 truly reveal dynamical dark energy? \NEW.
\item \textbf{arXiv:2511.22512} (November 2025, published March 2026) --- Robust DDE evidence from DESI DR2 + ACT + SPT + Planck. \NEW.
\item \textbf{Nature Astronomy (2025)} --- Dynamical dark energy in light of DESI DR2. \NEW.
\item \textbf{Euclid Prep.\ LXXXV} (A\&A, March 2026) --- DR1 higher-order WL. (continued)
\item \textbf{arXiv:2507.07629} --- ERO WL pipeline validation. (continued)
\item \textbf{Euclid QDR1} (March 2025) --- 30 TB data release. (continued)
\end{enumerate}


%=============================================================================
\part{Agent 13: Dark Sector}
\label{part:agent13}
%=============================================================================

\section{Mandate}

Neutrino mass hierarchy. Dark energy. Dark matter direct detection.

\section{\NEW: $\Sigma m_\nu$ Tension --- Landscape of Resolutions}

\begin{tcolorbox}[colback=yellow!5!white, colframe=yellow!60!black,
  title={$\Sigma m_\nu$ Tension: March 2026 Paper Explosion}]

The $\Sigma m_\nu$ tension has generated at least 5 new papers in March 2026 alone:

\begin{center}
\begin{tabular}{lp{5cm}p{4cm}}
\toprule
\textbf{Paper} & \textbf{Mechanism} & \textbf{Impact on DFD 0.061 eV} \\
\midrule
arXiv:2603.10787 & ACT DR6 + DESI DR2 in extended DE models & Safe in $w_0 w_a$CDM \\
arXiv:2603.13208 & Spatial curvature $\Omega_k$ & Tension drops to $1.17\sigma$ \\
arXiv:2506.22999 & Cosmic birefringence & Safe if birefringence real \\
arXiv:2601.04312 & Neutrino decays & Alternative explanation \\
arXiv:2603.03284 & Decaying dark matter & Background vs perturbations \\
\bottomrule
\end{tabular}
\end{center}

\textbf{DFD-preferred path}: The birefringence resolution (arXiv:2506.22999) is most interesting because it connects two DFD concerns simultaneously. The spatial curvature path (arXiv:2603.13208) is the most conservative.

\textbf{Assessment}: $\Sigma m_\nu$ remains \YELLOW{} but the risk is \textbf{lower} than in i46 because multiple credible resolution paths exist. DFD's 0.061 eV is not in existential danger.
\end{tcolorbox}

\section{JUNO: No New Data Since i46}

JUNO first results (59.1 days, arXiv:2511.14593) remain the latest. Normal ordering preferred at $2.2\sigma$. Next milestone: JUNO 1-year data ($\sim$2027).

\section{\NEW: Scalar NSI Risk to JUNO Mass Ordering}

\begin{tcolorbox}[colback=blue!5!white, colframe=blue!60!black,
  title={arXiv:2602.05564: Scalar Non-Standard Interactions Could Affect JUNO}]

A February 2026 paper (arXiv:2602.05564) shows that scalar non-standard interactions (NSI) could induce a resonance that complicates JUNO's mass ordering determination. The effect arises from interactions of neutrinos with the Earth's matter that modify the oscillation pattern.

\textbf{DFD assessment}: DFD does not predict scalar NSI of the type discussed. If scalar NSI were observed, it could be connected to DFD's $\psi$-field coupling to neutrinos, but this is speculative. The main concern is that scalar NSI could delay or complicate the mass ordering determination, extending the timeline for testing DFD's normal ordering prediction.

\textbf{Risk}: LOW. The paper identifies a theoretical possibility, not an observation. JUNO's baseline analysis assumes no NSI.
\end{tcolorbox}

\section{Agent 13 Assessment: Grade B}

Unchanged from i46. The $\Sigma m_\nu$ tension landscape analysis is the main contribution.

\section{New Literature (Agent 13)}

\begin{enumerate}
\item \textbf{arXiv:2603.10787} (March 2026) --- $\Sigma m_\nu$ with ACT DR6 + DESI DR2. \NEW.
\item \textbf{arXiv:2603.13208} (March 2026) --- Negative masses and spatial curvature. \NEW.
\item \textbf{arXiv:2602.05564} (February 2026) --- Scalar NSI risk to JUNO. \NEW.
\item \textbf{arXiv:2601.04312} (January 2026) --- Neutrino decays as $\Sigma m_\nu$ explanation. \NEW.
\item \textbf{arXiv:2603.03284} (March 2026) --- Decaying DM, background vs perturbations. \NEW.
\item \textbf{arXiv:2511.14593} (November 2025) --- JUNO first results. (continued)
\item \textbf{arXiv:2503.14744} (March 2025) --- DESI DR2 + Planck $\Sigma m_\nu$. (continued)
\end{enumerate}


%=============================================================================
\part{Agent 14: Numerics / SymBoltz / Phase 0}
\label{part:agent14}
%=============================================================================

\section{Mandate}

\textbf{CRITICAL}: Execute Phase 0A. Phase 0B framework.

\section{Phase 0A: EXECUTION STATUS}

\begin{tcolorbox}[colback=red!5!white, colframe=red!60!black,
  title={CRITICAL: Phase 0A --- Sixth Iteration Without Execution}]

\textbf{Plan quality}: Excellent. Fully specified since i46. 12 hours, 4 steps:
\begin{enumerate}
\item Hours 0--2: Install SymBoltz.jl v1.0.0. Verify $\Lambda$CDM baseline.
\item Hours 2--5: Implement distance-bias $D_L^{\text{DFD}} = D_L^{\Lambda\text{CDM}} \cdot e^{\Delta\psi(z)}$.
\item Hours 5--8: Compute $A_L^{\text{DFD}}$ and BAO distance ratios.
\item Hours 8--12: Generate Euclid pre-registration numbers.
\end{enumerate}

\textbf{Execution}: \textbf{ZERO}. The plan has been refined, extended, and praised for 6 iterations. Not a single line of code has been run.

\textbf{Root cause analysis}:
\begin{enumerate}
\item The iterative research cycle generates text, not code.
\item SymBoltz.jl installation and computation require a \textbf{dedicated computational session} with Julia environment setup, package installation, and numerical work.
\item This is fundamentally different from the text-generation tasks that the 17-agent programme excels at.
\end{enumerate}

\textbf{Recommendation}: Phase 0A must be \textbf{decoupled from the iterative cycle}. It requires a dedicated session with:
\begin{itemize}
\item Julia 1.10+ installed.
\item SymBoltz.jl cloned and built.
\item 12 uninterrupted hours of computational work.
\item A single focused agent, not 17 parallel agents.
\end{itemize}

\textbf{Deadline pressure}: Euclid pre-registration in 88 days. If Phase 0A is not executed by i49 (2 iterations from now), the programme will have failed on its single most important deliverable.
\end{tcolorbox}

\section{SymBoltz.jl: Published and Mature}

SymBoltz.jl v1.0.0 is now published in A\&A (March 2026). Key capabilities confirmed:
\begin{itemize}
\item Symbolic-numeric interface for custom models.
\item Implicit stiff ODE solvers (no approximation switching).
\item Automatic differentiation support.
\item Custom submodel support --- ideal for DFD's $\psi$-field extension.
\end{itemize}

\section{\NEW: gCAMB --- GPU-Accelerated Alternative}

\begin{tcolorbox}[colback=blue!5!white, colframe=blue!60!black,
  title={gCAMB (2026): GPU-Accelerated Boltzmann Solver}]

A new GPU-accelerated Boltzmann solver (gCAMB, published in A\&C, 2026) provides 10--100$\times$ speedup over CPU-based solvers. This is relevant for Phase 0A MCMC parameter estimation, though not for the initial background-only computation.

\textbf{DFD relevance}: For the Phase 0A Fisher forecast and eventual MCMC, GPU acceleration could be beneficial. However, SymBoltz.jl's differentiability (automatic differentiation) is more important for DFD's custom $\psi$-field model. gCAMB is a ``nice to have'', not a replacement for SymBoltz.jl.
\end{tcolorbox}

\section{Phase 0B: Framework --- Unchanged}

The five parameters ($m_\psi$, $c_s^2$, $\pi_\psi$, photon/baryon coupling, initial conditions) remain underived. The key literature (K-essence Boltzmann dynamics, $F(T)$ gravity in Boltzmann codes) remains identified. No progress on the derivation.

\section{Agent 14 Assessment: Grade C+}

Downgrade from B to C+ due to six iterations without execution. The plan quality is A-level. The execution is F-level. The average is C+.

\section{New Literature (Agent 14)}

\begin{enumerate}
\item \textbf{gCAMB} (A\&C, 2026) --- GPU-accelerated Boltzmann solver. \NEW.
\item \textbf{SymBoltz.jl} (A\&A, March 2026) --- Published. (continued; status: published)
\item \textbf{arXiv:2602.19576} (February 2026) --- K-essence Boltzmann dynamics. (continued)
\item \textbf{arXiv:2512.16404} (December 2025) --- $F(T)$ gravity in Boltzmann codes. (continued)
\end{enumerate}


%=============================================================================
\part{Observables Summary: i46 $\to$ i47}
%=============================================================================

\section{Changes from i46}

\begin{center}
\begin{tabular}{lll}
\toprule
\textbf{Item} & \textbf{i46 Status} & \textbf{i47 Status} \\
\midrule
Th-229 programme & On track & On track (consolidated) \\
Moriond 2026 & Fully digested & Fully digested (no change) \\
Wide binaries & Converging to GR & \textbf{New triple modelling paper: GR consistent} \\
Birefringence & CS assessment done & \textbf{$n\pi$ ambiguity + galaxy probe + $\Sigma m_\nu$ link} \\
$A_L$ & Phase 0A needed & Phase 0A needed (no change) \\
Euclid & 88-day countdown & 88-day countdown (qualitative pre-reg recommended) \\
$\Sigma m_\nu$ & FC tension noted & \textbf{Three resolution paths; risk reduced} \\
DESI DE & $2.8$--$4.2\sigma$ & \textbf{Robustness debate intensifies} \\
JUNO & First results & First results (no new data) \\
Phase 0A & Plan specified & \textbf{STILL UNEXECUTED --- 6 iterations} \\
Phase 0B & Framework outlined & Framework outlined (no change) \\
\bottomrule
\end{tabular}
\end{center}

\section{Observables Sector Assessment}

\begin{itemize}
\item \textbf{Strengths}: Laboratory programme on track. No BSM at LHC. Wide binaries converging to GR. Multiple $\Sigma m_\nu$ resolution paths.
\item \textbf{Weaknesses}: Phase 0A unexecuted for 6 iterations. $A_L$ blocked. Euclid pre-registration at risk.
\item \textbf{Opportunities}: DESI DE signal could be DFD $\psi$-field distance bias. Birefringence resolution of $\Sigma m_\nu$ tension could benefit DFD.
\item \textbf{Threats}: $n\pi$ birefringence ambiguity if angle is large. Phase 0A remaining unexecuted through i51.
\end{itemize}

\end{document}
